% Method. Your draft of 2026-10-25, mechanical fixes only --- see intro.tex.
% The detailed pipeline (configs, oracles, reward aggregation) stays in
% methods/overview.tex, which this section currently duplicates in part.
\section{Method}

We find that we can train for the objective of maximizing the ability of
fragments from ProteinMPNN\todo{sentence is incomplete --- maximizing their
ability to \emph{what}? recombine? fold when recombined?}. To do this we take a
protein of length $L$ and we split it up into $K$ fragments such that each
fragment $f_i$ is approximately $L/K$ amino acids long. Fragments are
purposely built to be equal length to make DNA assembly easier, and empirically
we do not find that it matters so much how you split up the protein (i.e.\
secondary structure, domains, etc.). We find that it is possible to find a
space of protein fragments that are highly combinable while also still diverse.
By treating the ProteinMPNN decoder as a policy conditioned on the encoded
backbone, we can do a short RL trajectory to align the decoder to the space of
highly compatible combinatorial fragments. For an encoded backbone we decode
$B$ sequences, then we split each into fragments, sample the fragments into a
larger pool of sequences $B \times 4$ such that each fragment is seen 4 times
in different combinations. When each sequence is folded a score is computed by
measuring RMSDs in the active site and confidence. We aggregate these into a
score for each fragment by averaging the score across all the sequences that
fragment was contained in. We provide rewards at the fragment level when we go
to update the policy. This means that each fragment in a given sequence will
receive a different reward. In practice there might be some issues here because
as the fragments are being decoded they will condition on other fragments as
well, but empirically this objective seems to be stable and optimizes with the
GRPO objective~\cite{shao2024deepseekmath}.

\plan{Three details from the implementation that this paragraph does not yet
state, all of which a referee will ask for.
(i)~\textbf{The pool contains unrecombined parents.} \texttt{build\_chimeras}
maps $(\rho, i)$ to parent $(i + \rho j) \bmod B$, so the $\rho = 0$ rotation
returns the parent itself --- 32 of every 128 chimeras at every $k$. Whether
that is deliberate (an anchor on the parental distribution) is worth one
sentence either way.
(ii)~\textbf{Only $k = 4$ was ever trained.} Everything in this paper about
granularity re-cuts $k = 4$ designs post hoc.
(iii)~\textbf{The fold is best-of-5.} That is the AME benchmark's own unit;
a single-sample run measures something the benchmark does not.}
